\chapter{JSON API, CORS, and browser client integration}

\noindent\textbf{Learning path: Fast path: API branch.} Read now for JSON/API clients or a separate frontend origin.\par\medskip

\rustbridgeblock{Typed scalars, structs/models, collections, \code{Result}, and reusable business functions are the same application core you use for HTML routes.}{JSON input mode, closed request schemas, explicit public projection, content negotiation, and CORS policy define the API boundary.}{Do not build a second business layer for JSON. Reuse the same typed logic and change only the representation and boundary contract.}


\invariantblock{Security invariant: CORS is not authorization}{CORS controls which browser origins may make or read cross-origin requests. It does not authenticate a caller, authorize an object, or replace the normal request lifecycle.}

\diagnosticblock{Security diagnostic}{SEC-DATA-004}{A model is being exposed without an explicit public projection, so internal fields could cross the response boundary unintentionally.}{Define and return the intended public projection instead of serializing the internal model shape directly.}

Velran does not force you to choose between classic server-rendered HTML and a JSON API. The same application can contain pages returning \code{Html}, actions processing forms, and typed JSON endpoints. This is useful both when a small JavaScript component needs some data and when a separate frontend application uses the Velran server as an API.

This chapter cleanly separates three operating models:

\begin{enumerate}
  \item \textbf{same-origin HTML + API}: the browser loads the page and calls the API from the same origin;
  \item \textbf{cross-origin public API}: another origin calls the API, but there is no session cookie;
  \item \textbf{cross-origin session auth}: a separate frontend origin uses cookie-based sessions, so CORS and CSRF policy are both required.
\end{enumerate}

The most important rule is: \textbf{CORS is neither authentication nor CSRF protection}. CORS is a browser-side access policy; authentication decides who the user is, while CSRF protection prevents an authenticated session from being abused for unauthorized state changes from another origin.

\section{The same lifecycle also applies to APIs}

JSON changes the representation, not the trust model. A state-changing API request should still be reasoned about in the same order:

\begin{lstlisting}
HTTP request
  -> parse: closed typed JSON schema
  -> validate: bounds/refinements
  -> authenticate: session or machine identity when required
  -> authorize: object/permission/tenant decision
  -> query / transaction
  -> response: explicit JSON projection or stable error envelope
\end{lstlisting}

CORS sits beside this lifecycle rather than replacing one of its stages. It answers whether a browser origin may make the response available to frontend code. It does not establish identity, object authority, CSRF safety, or database integrity. Keeping CORS outside the core authorization chain prevents a common design error: treating an allowed origin as an allowed user.

\section{The first typed JSON endpoint}

A JSON endpoint handler returns \code{Json}. \code{json(...)} converts a supported runtime value into a JSON representation.

\begin{lstlisting}[language=Velran]
#[page]
fn statusApi(ctx: PageContext) -> Result<Json, PageError> {
    let healthy = true;
    return Ok(json(healthy));
}

route statusApi GET "/api/status" public => statusApi;
\end{lstlisting}

A successful response is JSON, for example:

\Needspace{9\baselineskip}
\begin{lstlisting}
HTTP/1.1 200 OK
Content-Type: application/json; charset=utf-8

true
\end{lstlisting}

\code{json(value)} is not an arbitrary object serializer. Scalars can be encoded directly, but database/domain models cross a public JSON boundary through an explicit \code{expose(...)} projection. This prevents a newly added model field---especially a sensitive one---from silently appearing in an existing API. The response shape is therefore deliberate rather than a reflection of an arbitrary runtime object.

\section{JSON lists from the database}

Typed SQL and JSON responses connect directly. The query still returns an exact model shape, and the JSON layer represents that typed value.

\begin{lstlisting}[language=Velran]
model Product {
    id: i64
    name: String
    price: i64
}

#[query]
fn listProducts(db: Db, offset: i64)
    -> Result<Vec<Product>, DbError> sql {
    SELECT id, name, price
    FROM products
    ORDER BY id
    LIMIT 100 OFFSET :offset
}

#[page]
fn productsApi(ctx: PageContext, db: Db,
                    offset: i64)
    -> Result<Json, PageError> {
    let products = listProducts(db, offset)?;
    return Ok(json(expose(products, id, name, price)));
}

route productsApi GET "/api/products"
    query offset<i64>
    validate offset range 0 1000000
    public => productsApi;
\end{lstlisting}

The important point is not merely JSON encoding. The reusable list query owns its static row bound; the client supplies only a bounded offset, not SQL cardinality authority. Route input remains typed and validated, and the response uses an explicit public projection. A newly added model field therefore does not silently become part of the API.

\Needspace{17\baselineskip}
\section{Typed JSON POST body}

A route declares \code{json} input mode for a JSON request body.

\begin{lstlisting}[language=Velran]
#[action]
fn echoApi(ctx: ActionContext,
                  name: String,
                  age: i64,
                  active: bool)
    -> Result<Json, PageError> {
    return Ok(json(name));
}

route echoApi POST "/api/echo"
    json name<String> age<i64> active<bool>
    validate name length 1 100 age range 0 150
    auth user
    => echoApi;
\end{lstlisting}

A valid request is:

\begin{lstlisting}
POST /api/echo HTTP/1.1
Content-Type: application/json
Accept: application/json
X-CSRF-Token: <session CSRF token>

{"name":"Alice","age":42,"active":true}
\end{lstlisting}

The JSON input schema is \textbf{closed}. Unknown or missing keys, duplicate keys, wrong scalar types, and unsupported nested objects, arrays, floats, or \code{null} values are errors. These are request-contract violations and therefore return \code{400 Bad Request}. A missing or incorrect \code{Content-Type} on a JSON route returns \code{415 Unsupported Media Type}. Field and request-size limits apply just as they do to other request bodies.

The \code{json}, \code{form}, and \code{upload} body modes are mutually exclusive on the same POST route, and the compiler enforces this. JSON expects \code{application/json}. Simple forms use URL-encoded form data; uploads use multipart form data. Each endpoint should have one unambiguous request contract.

\section{Content negotiation and 406}

The client can use the \code{Accept} header to state what representation it can receive. For a JSON route, supported examples include:

\begin{lstlisting}
Accept: application/json
Accept: application/*
Accept: */*
\end{lstlisting}

If a handler would return JSON but the client accepts only \code{text/html}, the server returns \code{406 Not Acceptable}. This is preferable to silently returning a representation the client explicitly rejected.

\section{API errors}

On JSON routes, runtime-classified application errors use a simple common envelope: a stable \code{error} code, while detailed diagnosis remains in server-side logs. Clients should not parse the English text of error messages; they should use the HTTP status and stable error code.

Once a route is known to return JSON, errors use a JSON envelope, for example:

\begin{lstlisting}
HTTP/1.1 401 Unauthorized
Content-Type: application/json; charset=utf-8

{"error":"unauthorized"}
\end{lstlisting}

Typical error codes include:

\begin{itemize}
  \item \code{bad_request};
  \item \code{unauthorized};
  \item \code{forbidden};
  \item \code{csrf_failed};
  \item \code{unsupported_media_type};
  \item \code{database_unavailable};
  \item \code{resource_limit};
  \item \code{internal_error}.
\end{itemize}

This is especially important for browser clients. An authenticated JSON route returns \code{401} to an anonymous request rather than an HTML login redirect. The frontend can therefore decide based on status code without having to interpret HTML as an API error.

\section{Same-origin: the default and simplest model}

If HTML and API are served from the same scheme + host + port combination, they are same-origin from the browser's point of view.

\begin{lstlisting}
https://app.example/
https://app.example/api/products
\end{lstlisting}

No CORS configuration is normally needed. A simple client is:

\begin{lstlisting}
const response = await fetch("/api/products?limit=20&offset=0", {
  headers: { Accept: "application/json" }
});

if (!response.ok) {
  throw new Error(`HTTP ${response.status}`);
}

const products = await response.json();
\end{lstlisting}

Prefer same-origin unless a separate frontend origin is a real requirement; it removes a trust boundary and most CORS complexity.

\section{CORS: exact-origin allowlist}

Cross-origin browser access is denied by default. Trusted server configuration must enumerate exact origins; there is no wildcard origin.

\begin{lstlisting}
[web]
cors_origins = ["https://frontend.example"]
cors_allow_credentials = false
\end{lstlisting}

The CLI equivalent is:

\begin{lstlisting}
velran-server --cors-origin https://frontend.example ...
\end{lstlisting}

List every permitted origin explicitly; never reflect an arbitrary request \code{Origin}.

\section{Preflight}

Before some cross-origin requests, the browser sends an \code{OPTIONS} preflight. Velran allows the preflight only when:

\begin{itemize}
  \item the \code{Origin} is configured;
  \item the target GET/POST route exists;
  \item the requested request headers are supported.
\end{itemize}

Supported headers include \code{Content-Type}, \code{Accept}, and \code{X-CSRF-Token}; preflight remains part of route/origin policy, not a generic allow-all OPTIONS endpoint.

\section{Credentialed CORS with a session cookie}

When a frontend on a different origin uses Velran session authentication, credentialed CORS is required:

\begin{lstlisting}
[web]
cors_origins = ["https://frontend.example"]
cors_allow_credentials = true
\end{lstlisting}

In production this is allowed only over HTTPS. With this configuration the session cookie uses \code{SameSite=None; Secure}, because the browser must send it on cross-site requests.

The JavaScript fetch call also needs explicit credential mode:

\begin{lstlisting}
const response = await fetch("https://api.example/api/account", {
  credentials: "include",
  headers: { Accept: "application/json" }
});
\end{lstlisting}

\code{credentials: "include"} only allows matching cookies to be sent; it grants no application authority.

\section{CSRF on a JSON API}

With cookie-based sessions, the browser manages the session cookie automatically, so state-changing JSON requests need CSRF protection just like HTML forms.

For form POST, the token is a body field:

\begin{lstlisting}
_csrf=<token>
\end{lstlisting}

For JSON POST, use a header:

\begin{lstlisting}
X-CSRF-Token: <token>
\end{lstlisting}

For a separate frontend origin, an authenticated token endpoint can be defined:

\begin{lstlisting}[language=Velran]
#[page]
fn csrfApi(ctx: PageContext) -> Result<Json, PageError> {
    return Ok(json(csrfToken));
}

route csrfApi GET "/api/csrf" auth user => csrfApi;
\end{lstlisting}

The token belongs to the same session; CORS does not replace CSRF validation.

\section{Complete browser POST flow}

For a separate session-auth frontend, the browser sends the session cookie and session-bound CSRF token; the runtime verifies origin, authentication, CSRF, and typed input before the handler performs the business operation.

Example client:

\begin{lstlisting}
const csrf = await fetch("https://api.example/api/csrf", {
  credentials: "include",
  headers: { Accept: "application/json" }
}).then(async r => {
  if (!r.ok) throw new Error(`CSRF HTTP ${r.status}`);
  return r.json();
});

const result = await fetch("https://api.example/api/echo", {
  method: "POST",
  credentials: "include",
  headers: {
    "Content-Type": "application/json",
    Accept: "application/json",
    "X-CSRF-Token": csrf
  },
  body: JSON.stringify({ name: "Alice", age: 42, active: true })
}).then(async r => {
  if (!r.ok) throw new Error(`API HTTP ${r.status}`);
  return r.json();
});
\end{lstlisting}

\section{Reverse proxy and origin}

Behind Apache or Nginx, public scheme/host and trusted-proxy configuration must agree. The application must not decide what origin is safe based on arbitrary forwarded headers supplied by the client. The installation chapter explains the trusted-proxy boundary; the full \code{server.toml} key reference appears near the end of the book.

Always place the \textbf{frontend origin as seen by the browser} in the CORS allowlist, for example:

\begin{lstlisting}
https://frontend.example
\end{lstlisting}

not an internal proxy or container address.

\section{What not to do}

\begin{itemize}
  \item Do not enable or reflect CORS permissively; enumerate exact origins.
  \item Do not treat CORS as authentication or CSRF protection, and do not send session credentials over HTTP in production.
  \item Keep each POST route to one typed body mode and bound API input.
  \item JSON clients should handle API status codes such as 401/403, not HTML login redirects.
\end{itemize}

\section{Which architecture should you choose?}

\begin{center}
\begin{tabularx}{\textwidth}{@{}L{0.27\textwidth}Y@{}}
\toprule
\textbf{Situation} & \textbf{Recommended model} \\
\midrule
server-rendered page + small JS & same-origin HTML + JSON API \\
SPA on the same domain & same-origin API, no CORS \\
separate public frontend, no session & exact-origin CORS without credentials \\
separate frontend + session auth & exact-origin credentialed CORS + HTTPS + CSRF \\
server-to-server client & CORS is irrelevant; auth and network policy matter \\
\bottomrule
\end{tabularx}
\end{center}

\section{Runnable example}

The complete repository example is \path{examples/json-api/app.vrn}.

Start same-origin; add a separate frontend origin and CORS only when deployment requires it.

\section{Verified companion examples}
The repository companions are:
\begin{itemize}
  \item \path{examples/json-api/app.vrn} -- JSON API route flow;
  \item \path{examples/typed-query-form/main.vrn} -- typed query/form structs;
  \item \path{examples/json-typed-response/app.vrn} -- typed JSON response boundary.
\end{itemize}
Prefer these files when you need a complete request/response example rather than an isolated listing.

\section{Practice: separate API design from CORS design}
\paragraph{Exercise mode: supported.} Use the supported method defined in the preface.

Describe the JSON request and response contract without mentioning browsers. Then, separately, decide whether a browser on another origin needs permission to call it. This prevents CORS from being mistaken for authentication or API authorization.

\section{Common mistake: enabling permissive CORS by default}
CORS is a browser-origin policy, not a general network security mechanism. Prefer same-origin deployment when practical; when cross-origin access is required, enumerate exact origins and credential behavior deliberately.

\section{Running project checkpoint: Secure Notes}
Start from \path{examples/secure-notes/checkpoints/09-json-api/main.vrn} and its verified JSON companions. Expose a read-only JSON representation of the authenticated user's notes. Keep the same authorization rules as the HTML page. Add CORS only if you intentionally deploy a separate frontend origin; otherwise leave the API same-origin.
